                                 IMC2011, Blagoevgrad, Bulgaria
                                                 Day 2, July 31, 2011

                                           1
Problem 1. Let (an )∞                                                                         ∞
                    n=0 be a sequence with 2 < an < 1 for all n ≥ 0. Define the sequence (xn )n=0 by

                                                              an+1 + xn
                                    x 0 = a0 ,      xn+1 =                 (n ≥ 0).
                                                             1 + an+1 xn

What are the possible values of lim xn ? Can such a sequence diverge?
                                  n→∞
                                                                                             Johnson Olaleru, Lagos

Solution 1. We prove by induction that
                                                                    1
                                                   0 < 1 − xn <        .
                                                                  2n+1
Then we will have (1 − xn ) → 0 and therefore xn → 1.
     The case n = 0 is true since 12 < x0 = a0 < 1.
     Supposing that the induction hypothesis holds for n, from the recurrence relation we get
                                                    an+1 + xn     1 − an+1
                                1 − xn+1 = 1 −                 =             (1 − xn ).
                                                   1 + an+1 xn   1 + an+1 xn

By
                                                   1 − an+1     1 − 12   1
                                            0<                <        =
                                                  1 + an+1 xn   1+0      2
we obtain
                                                      1            1   1     1
                                   0 < 1 − xn+1 <       (1 − xn ) < · n+1 = n+2 .
                                                      2            2 2     2
     Hence, the sequence converges in all cases and xn → 1.
Solution 2. As is well-known,
                                                              tanh u + tanh v
                                          tanh(u + v) =
                                                             1 + tanh u tanh v
for all real numbers u and v.
    Setting un = ar tanh an we have xn = tanh(u0 + u1 + · · · + un ). Then u0 + u1 + · · · + un > (n + 1)ar tanh 12
and lim xn = lim tanh u = 1.
     n→∞        u→∞
Remark. If the condition an ∈ ( 21 , 1) is replaced by an ∈ (0, 1) then the sequence remains increasing and bounded,
but the limit can be less than 1.
Problem 2. An alien race has three genders: male, female, and emale. A married triple consists of three persons,
one from each gender, who all like each other. Any person is allowed to belong to at most one married triple. A
special feature of this race is that feelings are always mutual — if x likes y, then y likes x.
     The race is sending an expedition to colonize a planet. The expedition has n males, n females, and n emales.
It is known that every expedition member likes at least k persons of each of the two other genders. The problem
is to create as many married triples as possible to produce healthy offspring so the colony could grow and prosper.

  a) Show that if n is even and k = n2 , then it might be impossible to create even one married triple.

  b) Show that if k ≥ 3n
                       4 , then it is always possible to create n disjoint married triples, thus marrying all of the
     expedition members.

                                                                            Fedor Duzhin and Nick Gravin, Singapore




                                                             1
Solution. (a) Let M be the set of males, F the set of females, and E the set of emales. Consider the (tripartite)
graph G with vertices M ∪ F ∪ E and edges for likes. A 3-cycle is then a possible family. We’ll call G the graph
of likes.
    First, let k = n2 . Then n has to be even and we need to construct a graph of likes with no 3-cycles. We’ll do
the following: divide each of the sets M , F , and E into two equal parts and draw all edges between two parts as
shown below:
                                                        F               E

                                                                             M

                                            M


                                                        E
                                                                        F


Clearly, there is no 3-cycle.
     (b) First divide the the expedition into male-emale-female triples arbitrarily. Let the unhappiness of such a
subdivision be the number of pairs of aliens that belong to the same triple but don’t like each other. We shall show
that if unhappiness is positive, then the unhappiness can be decreased by a simple operation. It will follow that
after several steps the unhappiness will be reduced to zero, which will lead to the happy marriage of everybody.
     Assume that we have an emale which doesn’t like at least one member of its triple (the other cases are similar).
We perform the following operation: we swap this emale with another emale, so that each of these two emales will
like the members of their new triples. Thus the unhappiness related to this emales will decrease, and the other
pairs that contribute to the unhappiness remain unchanged, therefore the unhappiness will be decreased.
     So, it remains to prove that such an operation is always possible. Enumerate the triples with 1, 2, . . . , n and
denote by Ei , Fi , Mi the emale, female, and male members of the ith triple, respectively. Without loss of generality
we may assume that E1 doesn’t like either F1 or M1 or both. We have to find an index i > 1 such that Ei likes
the couple F1 , M1 and E1 likes the couple Fi , Mi ; then we can swap E1 and Ei .
     There are at most n/4 indices i for which E1 dislikes Fi and at most n/4 indices for which E1 dislikes Mi , so
there are no more than n/2 indices i for which E1 dislikes someone from the couple Mi , Fi , and the set of these
undesirable indexes includes 1. Similarly, there are no more than n/2 indices such that either M1 or F1 dislikes
Ei . Since both undesirable sets of indices have at most n/2 elements and both contain 1, their union doesn’t cover
all indices, so we have some i which satisfies all conditions. Therefore we can always perform the operation that
decreases unhappiness.
Solution 2 (for part b). Suppose that k ≥ 3n       4 and let’s show that it’s possible to marry all of the colonists.
First, we’ll prove that there exists a perfect matching between M and F . We need to check the condition of Hall’s
marriage theorem. In other words, for A ⊂ M , let B ⊂ F be the set of all vertices of F adjacent to at least one
vertex of A. Then we need to show that |A| ≤ |B|. Let us assume the contrary, that is |A| > |B|. Clearly, |B| ≥ k
if A is not empty. Let’s consider any f ∈ F \ B. Then f is not adjacent to any vertex in A, therefore, f has degree
in M not more than n − |A| < n − |B| ≤ n − k ≤ n4 , a contradiction.
     Let’s now construct a new bipartite graph, say H. The set of its vertices is P ∪ E, where P is the set of pairs
male–female from the perfect matching we just found. We will have an edge from (m, f ) = p ∈ P to e ∈ E for
each 3-cycle (m, f, e) of the graph G, where (m, f ) ∈ P and e ∈ E. Notice that the degree of each vertex of P in
H is then at least 2k − n.
     What remains is to show that H satisfies the condition of Hall’s marriage theorem and hence has a perfect
matching. Assume, on the contrary, that the following happens. There is A ⊂ P and B ⊂ E such that |A| = l,
|B| < l, and B is the set of all vertices of E adjacent to at least one vertex of A. Since the degree of each vertex
of P is at least 2k − n, we have 2k − n ≤ |B| < l. On the other hand, let e ∈ E \ B. Then for each pair
(m, f ) = p ∈ P , at most one of the pairs (e, m) and (e, f ) is joined by an edge and hence the degree of e in G is
at most |M \ A| + |F \ A| + |A| = 2(n − l) + l = 2n − l. But the degree of any vertex of G is 2k and thus we get
2k ≤ 2n − l, that is, l ≤ 2n − 2k.
     Finally, 2k − n < l ≤ 2n − 2k implies that k < 3n4 . This contradiction concludes the solution.
Problem 3. Determine the value of
                                 ∞                                                     
                                 X            1                 1                    1
                                       ln 1 +         · ln 1 +          · ln 1 +                .
                                              n                2n                  2n + 1
                                 n=1

                                                                                                    Gerhard Woeginger, Utrecht

                                                                2
Solution. Define f (n) = ln( n+1
                              n ) for n ≥ 1, and observe that f (2n)+f (2n+1) = f (n). The well-known inequality
ln(1 + x) ≤ x implies f (n) ≤ 1/n. Furthermore introduce
                                                  2n−1
                                                   X
                                         g(n) =          f 3 (k) < n f 3 (n) ≤ 1/n2 .
                                                  k=n

Then

                          g(n) − g(n + 1) = f 3 (n) − f 3 (2n) − f 3 (2n + 1)
                                              = (f (2n) + f (2n + 1))3 − f 3 (2n) − f 3 (2n + 1)
                                              = 3 (f (2n) + f (2n + 1)) f (2n) f (2n + 1)
                                              = 3 f (n) f (2n) f (2n + 1),
therefore
                   N                                       N
                   X                                 1 X                   1
                         f (n) f (2n) f (2n + 1) =       g(n) − g(n + 1) =   (g(1) − g(N + 1)) .
                                                     3                     3
                   n=1                                   n=1
Since g(N + 1) → 0 as N → ∞, the value of the considered sum hence is
                                   ∞
                                   X                                 1        1 3
                                         f (n) f (2n) f (2n + 1) =     g(1) =   ln (2).
                                                                     3        3
                                   n=1


                                                                                             f (k) − f (m)
Problem 4. Let f (x) be a polynomial with real coefficients of degree n. Suppose that                      is an integer
                                                                                                 k−m
for all integers 0 ≤ k < m ≤ n. Prove that a − b divides f (a) − f (b) for all pairs of distinct integers a and b.
                                                                                        Fedor Petrov, St. Petersburg

Solution 1. We need the following
Lemma. Denote the least common multiple of 1, 2, . . . , k by L(k), and define
                                                     
                                                       x
                                  hk (x) = L(k) ·              (k = 1, 2, . . .).
                                                       k
Then the polynomial hk (x) satisfies the condition, i.e. a − b divides hk (a) − hk (b) for all pairs of distinct integers
a, b.
Proof. It is known that
                                             X    k               
                                            a            a−b        b
                                                =                        .
                                            k              j      k−j
                                                         j=0

(This formula can be proved by comparing the coefficient of xk in (1 + x)a and (1 + x)a−b (1 + x)b .) From here we
get
                                          k                            k                        
                            a     b            X      a−b        b              X   L(k) a − b − 1        b
   hk (a) − hk (b) = L(K)     −        = L(K)                         = (a − b)                                .
                            k     k                    j      k−j                     j      j−1         k−j
                                                         j=1                            j=1

On the right-hand side all fractions L(k)
                                        j are integers, so the right-hand side is a multiple of (a, b). The lemma is
proved.
   Expand the polynomial f in the basis 1, x1 , x2 , . . . as
                                                 

                                                                            
                                                    x          x                x
                                  f (x) = A0 + A1       + A2      + · · · + An     .                            (1)
                                                    1          2                n
We prove by induction on j that Aj is a multiple of L(j) for 1 ≤ j ≤ n. (In particular, Aj is an integer for j ≥ 1.)
Assume that L(j) divides Aj for 1 ≤ j ≤ m − 1. Substituting m and some k ∈ {0, 1, . . . , m − 1} in (1),
                                                     m−1
                                  f (m) − f (k)   X Aj hj (m) − hj (k)    Am
                                                =        ·             +     .
                                     m−k            L(j)   m−k           m−k
                                                     j=1


                                                               3
                                                  Am
Since all other terms are integers, the last term m−k is also an integer. This holds for all 0 ≤ k < m, so Am is an
integer that is divisible by L(m).
   Hence, Aj is a multiple of L(j) for every 1 ≤ j ≤ n. By the lemma this implies the problem statement.
Solution 2. The statement of the problem follows immediately from the following claim, applied to the polynomial
g(x, y) = f (x)−f
              x−y
                  (y)
                      .
Claim. Let g(x, y) be a real polynomial of two variables with total degree less than n. Suppose that g(k, m) is an
integer whenever 0 ≤ k < m ≤ n are integers. Then g(k, m) is a integer for every pair k, m of integers.
Proof. Apply induction on n. If n = 1 then g is a constant. This constant can be read from g(0, 1) which is an
integer, so the claim is true.
   Now suppose that n ≥ 2 and the claim holds for n − 1. Consider the polynomials

                    g1 (x, y) = g(x + 1, y + 1) − g(x, y + 1)     and g2 (x, y) = g(x, y + 1) − g(x, y).                 (1)

For every pair 0 ≤ k < m ≤ n − 1 of integers, the numbers g(k, m), g(k, m + 1) and g(k + 1, m + 1) are all
integers, so g1 (k, m) and g2 (k, m) are integers, too. Moreover, in (1) the maximal degree terms of g cancel out, so
deg g1 , deg g2 < deg g. Hence, we can apply the induction hypothesis to the polynomials g1 and g2 and we thus
have g1 (k, m), g2 (k, m) ∈ Z for all k, m ∈ Z.
   In view of (1), for all k, m ∈ Z, we have that

 (a) g(0, 1) ∈ Z;

 (b) g(k, m) ∈ Z if and only if g(k + 1, m + 1) ∈ Z;

 (c) g(k, m) ∈ Z if and only if g(k, m + 1) ∈ Z.

For arbitrary integers k, m, apply (b) |k| times then apply (c) |m − k − 1| times as

                              g(k, m) ∈ Z ⇔ . . . ⇔ g(0, m − k) ∈ Z ⇔ . . . ⇔ g(0, 1) ∈ Z.

Hence, g(k, m) ∈ Z. The claim has been proved.

Problem 5. Let F = A0 A1 . . . An be a convex polygon in the plane. Define for all 1 ≤ k ≤ n − 1 the operation
fk which replaces F with a new polygon

                                          fk (F ) = A0 . . . Ak−1 A′k Ak+1 . . . An ,

where A′k is the point symmetric to Ak with respect to the perpendicular bisector of Ak−1 Ak+1 . Prove that
(f1 ◦ f2 ◦ . . . ◦ fn−1 )n (F ) = F . We suppose that all operations are well-defined on the polygons, to which they are
applied, i.e. results are convex polygons again. (A0 , A1 , . . . , An are the vertices of F in consecutive order.)
                                                                                        Mikhail Khristoforov, St. Petersburg

Solution. The operations fi are rational maps on the 2(n − 1)-dimensional phase space of coordinates of the
vertices A1 , . . . , An−1 . To show that (f1 ◦ f2 ◦ . . . ◦ fn−1 )n is the identity, it is sufficient to verify this on some
open set. For example, we can choose a neighborhood of the regular polygon, then all intermediate polygons in
the proof will be convex.
   Consider the operations fi . Notice that (i) fi ◦ fi = id and (ii) fi ◦ fj = fj ◦ fi for |i − j| ≥ 2. We also show
that (iii) (fi ◦ fi+1 )3 = id for 1 ≤ i ≤ n − 1.
   The operations fi and fi+1 change the order of side lengths by interchanging two consecutive sides; after
performing (fi ◦ fi+1 )3 , the side lengths are in the original order. Moreover, the sums of opposite angles in
the convex quadrilateral Ai−1 Ai Ai+1 Ai+2 are preserved in all operations. These quantities uniquely determine
the quadrilateral, because with fixed sides, both angles ∠A1 A2 A3 and ∠A1 A4 A3 decrease when A1 A3 increases.
Hence, property (iii) is proved.
   In the symmetric group Sn , the transpositions σi = (i, i + 1), which from a generator system, satisfy the same
properties (i–iii). It is well-known that Sn is the maximal group with n − 1 generators, satisfying (i–iii). In Sn we
have (σ1 ◦ σ2 ◦ . . . ◦ σn−1 )n = (1, 2, 3, . . . , n)n = id, so this implies (f1 ◦ f2 ◦ . . . ◦ fn−1 )n = id.



                                                              4
